%%%%%%%%%%%%%%%%%%%%%%%%%%%%%%%%%%%%%%%%%%%%%%%%%
\section{引言}
\label{sec:intro}

部分子分布函数（PDF）$f_{i/h}(x,\mu^2)$ 定义为：在因子化标度 $\mu$ 下探测，在快速运动的强子 $h$ 中找到一个夸克、反夸克或胶子（分别对应 $i=q,\bar{q},g$）、且其携带的强子动量分数介于 $x$ 与 $x+dx$ 之间的概率分布 \cite{Collins:1989gx}。它们是 QCD 中基本而重要的非微扰函数，定量刻画了强子与其内部夸克、胶子之间的关系，并在高能散射过程中把碰撞的强子（们）与短程 QCD 动力学联系起来，发挥着不可或缺的作用 \cite{Brambilla:2014jmp}。然而，PDF 包含长程尺度上的动力学、本质上是非微扰的，并且由含时算符定义，因此无论解析地计算还是从格点 QCD（LQCD）出发，从第一性原理计算 PDF 即便并非不可能，也是一大挑战。传统上，PDF 是在 QCD 因子化框架下通过对高能散射数据作 QCD 全局分析提取出来的 \cite{Dulat:2015mca,Martin:2009iq,Ball:2014uwa,Alekhin:2013nda,Ethier:2017zbq}。

最近，Ji 为动量为 $p_z$（比如沿 $z$ 方向）的强子引入了一组准 PDF（quasi-PDF），并论证当强子动量 $p_z$ 趋于无穷时它们等于相应的 PDF \cite{Ji:2013dva}。
在不取 $p_z\to\infty$ 极限的情况下，只要准 PDF 能够被乘法重正化，准 PDF 在 QCD 微扰论所有阶上都可以因子化为 PDF \cite{Ma:2014jla}。
与 PDF 类似，准 PDF 由两场关联量的强子矩阵元定义，两场之间有一条直线规范链接以保证规范不变性：
\begin{eqnarray}
\label{eq:ftq}
\tilde{f}_{q/p}(\tilde{x},\tilde{\mu}^2,p_z)
&\equiv &
\int \frac{d\xi_z}{2\pi} e^{i \tilde{x}p_z \xi_z}\langle h(p)| \overline{\psi}_q(\xi_z)\,\frac{\gamma_z}{2}
\nonumber\\
&\ & \times
\Phi_{n_z}^{(f)}(\{\xi_z,0\})\, \psi_q(0) | h(p) \rangle
\end{eqnarray}
为夸克准分布，其中 $\xi_0=\xi_\perp=0$；而
\begin{eqnarray}
\label{eq:ftg}
\tilde{f}_{g/p}(\tilde{x},\tilde{\mu}^2,p_z)
&\equiv &
\frac{1}{\tilde{x}p_z}
\int \frac{d\xi_z}{2\pi} e^{i \tilde{x} p_z \xi_z}\langle h(p)| F_{z}^{\ \nu}(\xi_z)\,
\nonumber\\
&\ & \times
\Phi_{n_z}^{(a)}(\{\xi_z,0\})\, F_{z\nu}(0) | h(p) \rangle
\end{eqnarray}
为胶子准分布，其中 $\nu$ 取遍横向方向。
在式~(\ref{eq:ftq}) 与 (\ref{eq:ftg}) 中，$\Phi_{n_z}^{(f,a)}(\{\xi_z,0\}) = \text{exp}[-ig\int_0^{\xi_z} d\eta_z\, A^{(f,a)}_z(\eta_z)]$ 是规范链接，``$f$''与``$a$''分别表示基础表示和伴随表示。
与 PDF 不同，准 PDF 的两场关联量定义在光锥之外且具有等时间隔，这使得在 LQCD 中计算准 PDF 成为可能 \cite{Lin:2014zya,Alexandrou:2015rja,Chen:2016utp,Zhang:2017bzy,Orginos:2017kos}。
然而，由于 LQCD 计算中的强子动量实际上受格距的限制，$p_z\to\infty$ 的极限在 LQCD 计算中难以达到，控制有限 $p_z$ 带来的修正是一个挑战 \cite{Ji:2014gla,Ji:2017rah}。尽管如此，凭借利用 LQCD 从 QCD 第一性原理计算 PDF 的潜力以及随之而来的挑战，准 PDF 的概念在微扰 QCD（PQCD）与 LQCD 两个学界都引发了大量的兴趣与活动 \cite{Ji:2014hxa,Gamberg:2014zwa,Jia:2015pxx,Li:2016amo,Bacchetta:2016zjm,Radyushkin:2016hsy,Radyushkin:2017gjd,Gbedo:2017eyp,Radyushkin:2017ffo,Carlson:2017gpk,Briceno:2017cpo,Nam:2017gzm,Xiong:2017jtn,Radyushkin:2017cyf,Constantinou:2017sej,Yoon:2017qzo,Rossi:2017muf}。除了 $p_z$ 取值范围有限带来的修正之外，从 LQCD 计算出的准 PDF 导出 PDF 的关键有两点：(1) 能够非微扰地重正化准 PDF 的所有紫外（UV）发散；(2) 保证重正化后的准 PDF 与 PDF 共享同样的共线（CO）发散。

本文的两位作者曾在文献 \cite{Ma:2014jla} 中证明：只要准 PDF 能够被乘法重正化，准 PDF 与相应的 PDF 在 QCD 微扰论所有阶上都共享同样的领头对数的共线（CO）微扰发散。基于这一事实，两位作者又在文献 \cite{Ma:2014jla} 中提出：不必通过取强子动量 $p_z\to\infty$ 从准 PDF 得到 PDF，而是采用 PQCD 因子化方法，从 LQCD 计算出的准 PDF 或其他 LQCD 可计算的单强子矩阵元中提取 PDF，只要这些矩阵元能够以微扰可计算的系数因子化为所需的 PDF。在这种 PQCD 因子化途径中，硬标度为 $\tilde{x}p_z \sim 1/\xi_z \sim \text{GeV}$ 量级，修正按硬标度的幂次压低，并由高扭度算符的强子矩阵元刻画，一如实验可测且微扰可因子化的强子截面经由 PQCD 因子化与 PDF 相联系的方式。

因此，理解定义准 PDF 的算符的可重正性，是从 LQCD 计算出的准 PDF 可靠地提取 PDF 的最关键挑战。由于准 PDF 不是由扭度2算符定义的，PDF 与准 PDF 具有不同的微扰紫外行为：准 PDF 不是像 PDF 那样具有对数的微扰紫外发散，而是具有幂次紫外发散 \cite{Ma:2014jla,Xiong:2013bka}。虽然准 PDF 的 LQCD 计算天然被格距 $a$ 正规化，但 $1/a$ 的微扰幂次发散使得格点结果难以取连续极限，从而难以提取正确的 PDF \cite{Ishikawa:2016znu,Chen:2016fxx,Monahan:2016bvm,Alexandrou:2017huk, Chen:2017mzz}。在紫外重正化下，定义准 PDF 的算符也可能与其他算符混合，例如夸克 PDF 可以与胶子 PDF 混合。因此，非常重要的问题是：不仅要弄清定义准 PDF 的算符是否可重正，还要弄清重正化下的算符混合（如果存在的话）能否限于一个封闭的算符集合。本文将处理与定义准 PDF 的算符的可重正性有关的问题。


本文其余部分安排如下。
在第 \ref{sec:def} 节我们将证明，由式~\eqref{eq:ftq} 与 \eqref{eq:ftg} 定义的准 PDF 具有糟糕的短距离行为，而坐标空间准 PDF 是提取 PDF 更好的候选者。我们还在坐标空间中定义准 PDF，并解释其重正化为何困难。随后在第 \ref{sec:oneloop} 节研究坐标空间准 PDF 的一圈展开。利用第 \ref{sec:powercounting} 节导出的紫外幂次计数，我们辨认出准 PDF 的所有微扰紫外发散区域。我们发现，一旦扣除子发散，所有紫外发散都来自各圈动量同时变大的积分区域，这与 PDF 的行为显著不同。
基于这些观察与发现，我们在第 \ref{sec:renormalization} 节证明坐标空间准 PDF 可以被乘法重正化。最重要的是，我们发现准 PDF 在重正化标度跑动下不与其他任何量混合，这完成了我们对坐标空间准 PDF 可重正性的证明。最后在第 \ref{sec:summary} 节给出总结。
